\documentclass[conference]{IEEEtran}

% Packages
\usepackage{graphicx}
\usepackage{amsmath}
\usepackage{amssymb}
\usepackage{booktabs}
\usepackage{multirow}
\usepackage{hyperref}
\usepackage{xcolor}
\usepackage{listings}

% Hyperref setup
\hypersetup{
    colorlinks=true,
    linkcolor=blue,
    citecolor=blue,
    urlcolor=blue
}

% Code listing setup
\lstset{
    basicstyle=\ttfamily\small,
    breaklines=true,
    frame=single,
    numbers=left,
    numberstyle=\tiny,
    captionpos=b
}

% Custom commands
\newcommand{\sharpe}{\text{Sharpe}}
\newcommand{\tsmom}{\text{TSMOM}}
\newcommand{\macd}{\text{MACD}}
\newcommand{\rsi}{\text{RSI}}


\begin{document}

\title{Evaluating Time Series Momentum in Multi-Agent Trading Systems: A Comparative Study with Technical Indicators}

\author{
\IEEEauthorblockN{AutoGen-Trader Research Team}
\IEEEauthorblockA{
AutoGen-Trader Project\\
GitHub: github.com/iAmGiG/AutoGen-Trader\\
}
}

\maketitle

\begin{abstract}
Time series momentum (TSMOM) has demonstrated consistent performance across asset classes with Sharpe ratios ranging from 0.48 to 0.79 in academic research (Moskowitz et al., 2012). However, its integration with technical indicator-based trading systems remains underexplored. This paper evaluates TSMOM as a complementary signal generator within a multi-agent trading architecture, comparing its performance against MACD-only and MACD+RSI dual voting strategies. Using a custom backtesting framework validated on 10 years of equity data (2016-2024), we test whether 12-month TSMOM signals achieve Sharpe > 0.6 with correlation < 0.5 to existing technical strategies. Our findings inform the design of triple voting consensus mechanisms and provide empirical validation of TSMOM in automated trading contexts.
\end{abstract}

\begin{IEEEkeywords}
Time series momentum, algorithmic trading, multi-agent systems, backtesting, technical analysis, voting strategies
\end{IEEEkeywords}

\section{Introduction}

Momentum strategies have been a cornerstone of quantitative finance since the 1990s \cite{jegadeesh1993returns}, with time series momentum (TSMOM) emerging as a particularly robust approach across diverse asset classes \cite{moskowitz2012time}. Unlike cross-sectional momentum, which compares assets relative to each other, TSMOM generates trading signals based solely on an asset's own historical price trajectory. This property makes TSMOM particularly attractive for automated trading systems operating on individual equities.

\subsection{Motivation}

The AutoGen-Trader project implements a multi-agent trading architecture where specialized agents generate independent signals that are combined through voting mechanisms. The current production system employs dual voting between MACD (Moving Average Convergence Divergence) and RSI (Relative Strength Index) indicators, achieving a Sharpe ratio of 0.856 on AAPL during 2024 validation. However, both MACD and RSI are technical indicators operating on similar timeframes (days to weeks), potentially exhibiting high correlation during regime shifts.

TSMOM operates on longer timeframes (typically 12 months), offering the potential for:
\begin{itemize}
    \item \textbf{Diversification}: Uncorrelated signals from different temporal perspectives
    \item \textbf{Regime Capture}: Longer lookback captures macro trends missed by short-term indicators
    \item \textbf{Consensus Robustness}: Third voter reduces false signals through triple voting
\end{itemize}

\subsection{Research Questions}

This paper addresses three key questions:

\begin{enumerate}
    \item \textbf{RQ1}: Does 12-month TSMOM achieve Sharpe > 0.6 as a standalone strategy on individual equities?
    \item \textbf{RQ2}: Is TSMOM correlation with MACD+RSI voting < 0.5, providing true diversification?
    \item \textbf{RQ3}: Does triple voting (MACD + RSI + TSMOM) improve risk-adjusted returns over dual voting?
\end{enumerate}

\subsection{Contributions}

Our work makes the following contributions:

\begin{itemize}
    \item Empirical validation of TSMOM on individual equities (2016-2024)
    \item Comparative analysis of TSMOM vs technical indicator strategies
    \item Design and validation of triple voting consensus mechanisms
    \item Open-source backtesting framework for multi-agent trading research
\end{itemize}

\subsection{Paper Organization}

Section II reviews related work on momentum strategies and multi-agent trading. Section III describes our methodology, including the backtesting framework and signal generation algorithms. Section IV details the experimental setup and validation approach. Section V presents results from TSMOM standalone testing and comparative analysis. Section VI discusses implications for multi-agent integration. Section VII concludes with contributions and future work.

\input{02_Related_work}
\input{03_Methodology}
\input{04_Experimental_setup}
\input{05_Results}
\input{06_Discussion}
\input{07_Conclusion}

\section*{Acknowledgments}
This research was conducted as part of the AutoGen-Trader project, an open-source multi-agent trading platform built on Microsoft AutoGen framework. We acknowledge the academic foundation provided by Moskowitz, Ooi, and Pedersen's seminal work on time series momentum.

\bibliographystyle{IEEEtran}
\bibliography{references}

\end{document}
